\chapter{Final Remarks and Future Work}\label{chap:final-remarks-and-future-work}\label{chap:conclusions}

This chapter summarizes the main contributions of this thesis, discusses the
key findings obtained from the computational experiments, and outlines
directions for future research.

\section{Summary of Contributions}

This thesis introduced and investigated the \gls{cbrp}, a new combinatorial
optimization problem that models the selection of city blocks for insecticide
spraying and the design of vehicle routes to service them. The \gls{cbrp}
combines features of both Vehicle Routing Problems and Arc Routing Problems,
since each city block must be completely encircled by the spraying vehicle and
only a subset of blocks can be serviced within a given time budget. Although
motivated by Dengue vector control, the \gls{cbrp} framework is applicable to a
variety of urban logistics problems, including waste collection, postal
delivery, and utility meter reading.

The contributions developed in this thesis span four interconnected research
axes: deterministic optimization, epidemiological simulation, stochastic
optimization, and simulation-optimization integration. The main contributions
along each axis are summarized below.

\paragraph{Deterministic optimization.}
The \gls{cbrp} was formally defined on a planar graph derived from real city
maps. Two families of \gls{milp} formulations were proposed: a path-based
family, in which no vertex is visited more than once, and a walk-based family,
in which vertices may be revisited. For each family, both an exponential
subtour-elimination formulation and a compact \gls{mtz} variant were developed.
Preprocessing techniques exploiting the planar structure of the underlying graph
were introduced, reducing the number of vertices by approximately 29\%, arcs by
27\%, and blocks by 65.5\% on average. To provide dual bounds without requiring
a general-purpose solver, four \gls{lr} variants were derived, decomposing the
problem into shortest-path and knapsack subproblems. On the primal side, a
greedy constructive heuristic combining binary search with knapsack-based block
selection and a block-connection routing heuristic based on shortest paths in a
\gls{dag} was designed. A deterministic benchmark comprising 39 instances
derived from real Dengue notification data from Alto Santo and Limoeiro, in
Cear\'{a}, Brazil, was constructed to support the experimental evaluation.

\paragraph{Epidemiological simulation.}
A \gls{mabs} grounded in compartmental epidemiological theory (SIR for humans,
SEI for mosquitoes) was developed to simulate the spatial-temporal progression
of Dengue cases at the city-block level. The simulator, implemented on the GAMA
platform, models the interactions among mosquitoes, humans, and breeding sites
over 12-hour cycles using real geographic data extracted from OpenStreetMap. The
simulator was calibrated and validated against historical notification data from
Alto Santo and Limoeiro, achieving strong Pearson correlation coefficients
(up to $r = 0.93$ for Alto Santo and $r = 0.97$ for Limoeiro) and
block-level accuracy of 86--94\% of blocks within the simulated endemic
channel. The simulation also provided insights into the effectiveness of
emergency fogging interventions.

\paragraph{Stochastic optimization.}
The \gls{cbrp} was extended to its stochastic counterpart, the \gls{scbrp},
formulated as a two-stage stochastic program. In the first stage, blocks are
selected for spraying under current conditions; in the second stage, independent
routes are computed for a set of future scenarios with a reduction factor
capturing the effectiveness of first-stage nebulization. Path-based and
walk-based \gls{milp} formulations, each with \gls{mtz} compact variants, were
proposed for the \gls{scbrp}. Tailored heuristics were developed, including an
initial solution heuristic based on the deterministic greedy approach, a local
search with multiple neighborhood structures using weak and moderate delta
evaluation strategies, and a \gls{sa} metaheuristic that balances
intensification and diversification through a temperature-controlled acceptance
mechanism. A stochastic benchmark comprising 36 instances, each with 50
scenarios generated by the \gls{mabs}, was constructed.

\paragraph{Simulation-optimization integration.}
A \gls{sh} framework was designed to bridge the simulation and optimization
components. The framework operates as an iterative process in which the
\gls{mabs} generates new scenarios that are progressively loaded into a
multi-start heuristic. The multi-start procedure evaluates six
profit-assignment strategies that combine deterministic case data with
simulation-derived statistics (accumulated cases and incidence), using both
unperturbed baselines and stochastic perturbations to diversify the search.
Each candidate solution is constructed via the greedy heuristic, refined through
a three-phase local search (route improvement, block insertion, and block swap),
and evaluated using a stochastic objective that accounts for expected future
demand.

\section{Key Findings}

The computational experiments conducted in this thesis yielded several
important findings.

\paragraph{Deterministic formulations.}
Among the path-based formulations, Path-CBRP-MTZ-Prep (the compact \gls{mtz}
variant with preprocessing) achieved the best overall performance, with an
average optimality gap of 2.97\% and 18 out of 39 instances solved to
optimality within a one-hour time limit. Path-based formulations consistently
outperformed walk-based ones: the best path-based configuration achieved an
average gap of 2.52\% and proved optimality in 19 instances, compared to an
average gap of 7.75\% and only 2 instances solved to optimality for the best
walk-based configuration. The preprocessing techniques proved beneficial for
most formulations, particularly the \gls{mtz} variants, while the fractional
separation cuts resulted in significantly worse performance (gaps approximately
six times higher than other path variants).

\paragraph{Lagrangean relaxation.}
The \gls{lr} variants provided competitive dual bounds: the best variant
(LR-RCSP-KN-HR-BM) outperformed the linear relaxation in 36 out of 39
instances while requiring less than 90 seconds of computation and no
general-purpose solver. The inclusion of a greedy heuristic for primal bounds
and the barrier method initialization for dual bounds were identified as the
most impactful algorithmic components.

\paragraph{Epidemiological simulation.}
The \gls{mabs} demonstrated the ability to reproduce the temporal dynamics of
Dengue spread at both the city and block levels. For Alto Santo, the simulator
achieved Pearson correlation coefficients ranging from 0.77 to 0.93 with mean
absolute errors between 0.94 and 8.40 weekly cases. For Limoeiro, the results
varied by period: strong correlations (up to $r = 0.97$) were obtained for
periods with higher incidence, while periods with irregular, low-incidence
patterns proved more challenging. At the block level, 86--94\% of blocks in
Alto Santo and 85--91\% in Limoeiro fell within the simulated endemic channel,
indicating that the simulator captures the spatial distribution of cases with
reasonable accuracy. Sensitivity analysis revealed that adult mosquito mortality
had the clearest monotonic effect on simulation outcomes, while aquatic
mortality and oviposition exhibited non-monotonic behavior.

\section{Future Work}

Several research directions emerge from this thesis.

\paragraph{Computational experiments for stochastic methodologies.}
A natural next step is the comprehensive computational evaluation of the
stochastic formulations, the \gls{sa} metaheuristic, and the \gls{sh}
framework on the stochastic benchmark. These experiments will provide
insights into the value of the stochastic solution, the expected value of
perfect information, and the relative performance of exact versus heuristic
approaches under uncertainty.

\paragraph{Multi-vehicle and multi-period extensions.}
The current \gls{cbrp} formulation considers a single vehicle with a fixed
time budget. Extending the problem to multiple vehicles operating simultaneously
and to multi-period planning horizons would better reflect the operational
reality of vector control campaigns, where fleets of spraying vehicles are
deployed over several consecutive days.

\paragraph{Dynamic and real-time routing.}
Integrating real-time surveillance data and weather information into the
routing decisions would enable dynamic re-routing as new Dengue cases are
reported or environmental conditions change. This extension could leverage the
simulation component to continuously update scenario estimates and trigger
route adjustments.

\paragraph{Enhanced simulation fidelity.}
The \gls{mabs} could be improved by incorporating additional environmental
factors such as rainfall patterns, temperature variations, and land-use data.
Coupling the simulation with epidemiological surveillance systems would allow
for automatic calibration and validation against real-time notification data.
Moreover, extending the simulation to model multiple Dengue serotypes and
co-circulating arboviruses (Zika, Chikungunya) would increase its applicability.

\paragraph{Algorithmic improvements.}
On the optimization side, developing branch-and-cut algorithms with tailored
valid inequalities for the \gls{scbrp} could improve the quality of exact
solutions. Investigating decomposition methods such as Benders decomposition or
column generation for the stochastic formulations may help address scalability
to larger instances. Additionally, exploring machine learning techniques to
predict block-level Dengue incidence and to warm-start the optimization
algorithms represents a promising research direction.

\paragraph{Practical deployment.}
Developing a decision support tool that integrates the optimization,
simulation, and data visualization components into a user-friendly platform
for health departments would facilitate the practical adoption of the
methodologies proposed in this thesis. Such a tool could assist public health
authorities in planning insecticide spraying campaigns, evaluating the expected
impact of different intervention strategies, and monitoring the progression of
Dengue outbreaks in real time.
